\pagestyle{plain}
\chapter*{Resumen}
En un mundo globalizado donde el auge de las tecnologías inalámbricas aumenta cada día, así como el número de dispositivos conectados a internet, se hace imprescindible conocer el estado de la seguridad de los protocolos de comunicación inalámbrica y los sistemas de automatización de hogares. Estos se ayudan de un amplio abanico de dispositivos conocidos como el internet de las cosas (IoT, del inglés Internet of Things).
\\\\
Dos de las tecnologías de IoT que han tenido un mayor crecimiento durante los últimos cinco años, han sido Thread y Zigbee. Thread es un protocolo basado en IPv6 diseñado para la comunicación entre múltiples dispositivos IoT conectados, mientras que Zigbee es un protocolo de comunicación inalámbrica para la radiodifusión digital de datos, diseñado para permitir el control y la monitorización de dispositivos. Estos protocolos inalámbricos están basados en el estándar IEEE 802.15.4, destacando en entornos de IoT domésticos e industriales debido a su bajo coste de implementación y consumo de energía.
\\\\
En este proyecto se ha realizado un estudio del estado del arte del estándar IEEE 802.15.4, que especifica las capas inferiores de la pila OSI (capa física [PHY] y capa de acceso al medio [MAC]), y de los protocolos Thread y Zigbee, haciendo énfasis en los mecanismos de seguridad utilizados para garantizar la confidencialidad de las comunicaciones entre los dispositivos IoT que utilizan estas tecnologías.
\\\\
Mediante la creación de un laboratorio de pruebas, se ha llevado ha cabo una auditoría de seguridad en una red industrial Zigbee y un análisis de las vulnerabilidades existentes en este tipo de redes, cuando se aplica en ellas una mala o nula configuración centrada en la seguridad de las comunicaciones y de la información que es enviada. Además, al realizar pruebas sobre la red, se descubrió y publicó posteriormente en Exploit-DB una vulnerabilidad desconocida hasta el momento, que tiene lugar en los dispositivos Zigbee del fabricante Ksix, cuya explotación permite falsificar mensajes entre dispositivos de forma remota eludiendo el mecanismo de protección ante repetición del protocolo.
\\\\
Finalmente, se han presentado un conjunto de medidas preventivas y/o correctivas necesarias para mitigar o eliminar los riesgos de sufrir ataques en redes IoT. Y se han establecido las conclusiones sobre el estado de la seguridad en las redes IoT y en las tecnologías analizadas: IEEE 802.15.4, Thread y Zigbee.
